\section*{ID6671: Optical Core: Aperture → Bessel → Airy Pattern: F\#: F\# = DfracFD}
\paragraph{Metadata.}
PiID: 9069963955245 | RTEID: RTE:P25904.9745:T4.9843 | UUID: 90d2b6f8-1f73-53e7-8526-a7d52f182abc

\paragraph{Canonical Statement.}
Relate to f-number \textbackslash{}(f\textbackslash{}\#\textbackslash{}): since \textbackslash{}(f\textbackslash{}\# = \textbackslash{}dfrac\{F\}\{D\}\textbackslash{}) (focal length \textbackslash{}(F\textbackslash{})), the *linear* radius of the Airy disk in the image plane is

\paragraph{Canonical Equation.}
\[
f\# = \dfrac{F}{D}
\]

\paragraph{Dictionary.}
Relate to f-number \textbackslash{}(f\textbackslash{} \textbackslash{}): since \textbackslash{}(f\textbackslash{} = FD\textbackslash{}) (focal length \textbackslash{}(F\textbackslash{})), the linear radius of the Airy disk in the image plane is (f\#)

\paragraph{Encyclopedia.}
Optical Core: Aperture → Bessel → Airy Pattern: F\#: F\# = DfracFD is an algebraic definition, written f\# = dfracFD. It is an algebraic definition expressing the quantity directly in terms of the variables on its right-hand side. It is built from F, D, defining f\#. Within the Trawin Zero Point registry it serves as one element of the framework's mathematical narrative.
